% !TeX encoding = UTF-8
\SetChapterInfo{subtitle={从紧性判别到 Fredholm 择一},keywords={全有界; 核算子; Rellich; Schauder; Fredholm; Dirichlet; 核空间},prereq={Ascoli 定理; Hahn--Banach; Sobolev 空间; 闭值域; Riesz 投影},goals={证明紧算子的谱结构; 识别相容条件; 用紧嵌入处理椭圆边值问题}}
\chapter{紧算子}\label{ch:yosida-10}
\ChapterSource{按 \cite{yosida-zh,yosida-en} 第十章六节及附录组织，2026 年 10 月 27 日补写为学习讲义。直接查看英文印刷 274--293 页全部扫描图，以及中译本印刷 287、290、294--296、299、302、306 页。两版参考范围分别为 287--306、274--293；英文第 274 页上部仍属第九章。查阅范围与未展开结果详见本章 COVERAGE。}

有界集在无限维空间中未必有收敛子列，紧算子却把每个有界序列送到具有收敛子列的序列。本章利用这一性质证明：紧算子在非零谱点附近保留有限维线性代数的结构，方程的障碍由有限多个对偶条件描述。应用中，椭圆方程的估计先给出有界解算子，Sobolev 紧嵌入再把它变为紧算子。

除另有说明，$X,Y$ 为实或复 Banach 空间，讨论谱时取复数域；$B_X$ 为闭单位球，$T'$ 为线性对偶算子。Hilbert 内积对第一变量线性，Hilbert 伴随写作 $T^*$。本章的核算子指 nuclear operator，不能与泛指由积分核定义的算子混同。

\Needspace{10\baselineskip}
\section{B-空间中的紧集}\label{sec:yosida-10-01}
\noindent\textit{原书起始页：中译本 287；英文版 274。}

\begin{definition}[全有界与相对紧]
集合 $M\subset X$ 称为全有界，若对每个 $\varepsilon>0$ 可用有限多个半径 $\varepsilon$ 的球覆盖；称为相对紧，若其闭包紧。在 Banach 空间中，两者等价。闭且有界本身不够，例如 $\ell^2$ 的单位向量满足 $\|e_n-e_m\|=\sqrt2$。
\end{definition}
\begin{proof}
紧性蕴含有限网。反之，逐次用半径 $2^{-j}$ 的有限覆盖作对角选取，得到 Cauchy 子列；完备性保证其收敛。因此闭包序列紧，因而紧。
\end{proof}

\begin{theorem}[Fréchet--Kolmogorov 判别]\label{thm:y10-kolmogorov}
设 $1\le p<\infty$。$M\subset L^p(\R^d)$ 相对紧，当且仅当以下三项成立：
\begin{align*}
 &\sup_{f\in M}\|f\|_p<\infty,\\
 &\lim_{h\to0}\sup_{f\in M}\|f(\cdot+h)-f\|_p=0,\\
 &\lim_{R\to\infty}\sup_{f\in M}\|1_{\{|x|>R\}}f\|_p=0.
\end{align*}
第二项排除越来越快的振荡，第三项排除质量逃向无穷远。
\end{theorem}
\begin{proof}
必要性：紧闭包有界。对单个 $f\in L^p$，平移连续性由 $C_c^\infty$ 的稠密性和移位等距性得到，尾部范数则由积分的绝对连续性趋于零。在 $M$ 中取有限 $\varepsilon$ 网，将这两个性质从有限个网点传给全体元素；平移差的误差至多 $2\varepsilon$，尾部误差至多 $\varepsilon$，故为一致极限。

充分性：取非负 $\rho\in C_c^\infty(B(0,1))$，$\int\rho=1$，令 $\rho_\delta(x)=\delta^{-d}\rho(x/\delta)$。Minkowski 不等式给出
\[
 \sup_{f\in M}\|\rho_\delta*f-f\|_p
 \le\sup_{|h|\le\delta,\,f\in M}\|f(\cdot-h)-f\|_p\longrightarrow0.
\]
固定 $\delta$，由 Hölder 不等式，函数 $g_f=\rho_\delta*f$ 及其一阶导数在全空间一致有界，常数可依赖 $\delta$。这对 $p=1$ 同样成立，因为卷积核及其导数有界。故 $\{g_f\}$ 在每个闭球上一致有界且等度连续。Ascoli 定理使其限制在该球上的集合在一致范数、从而在 $L^p$ 范数下全有界。

又当 $R>\delta$ 时，由卷积核的支集和 Minkowski 不等式，
\[
 \|1_{\{|x|>R\}}g_f\|_p
 \le\|1_{\{|x|>R-\delta\}}f\|_p.
\]
先选 $\delta$ 使平滑误差小，再选 $R$ 使尾部一致小，最后选球上的有限网并将网点在球外延零，得到 $M$ 的任意精度有限网。$L^p$ 完备，故 $M$ 相对紧。
\end{proof}

\begin{example}[紧性的两种损失]
固定非零 $f\in C_c^\infty(\R)$，$f_n(x)=f(x-n)$ 的平移连续性对 $n$ 一致，但尾部条件失败。另一方面，在 $L^2(\R)$ 中取 $g_n(x)=1_{[0,2\pi]}(x)e^{inx}$，它们有共同紧支集且范数相同，却彼此正交，故不相对紧；此时失败的是一致平移连续性。
\end{example}

\begin{exercise}[紧闭包的凸化]\label{ex:y10-convex}
证明 Banach 空间中全有界集的凸包仍全有界，其闭凸包因而紧。
\end{exercise}
\begin{solution}
取有限 $\varepsilon$ 网 $F$。每个有限凸组合与 $\operatorname{conv}F$ 中某点的距离至多 $\varepsilon$。$\operatorname{conv}F$ 位于有限维子空间且有界，故有有限 $\varepsilon$ 网；合成即得 $2\varepsilon$ 网。闭包的紧性使用环境空间完备性。
\end{solution}

\begin{exercise}[去掉尾部条件]\label{ex:y10-tail}
对上述 $f_n$，证明一致平移连续性，并用支集分离说明没有收敛子列。
\end{exercise}
\begin{solution}
$\|f_n(\cdot+h)-f_n\|_p=\|f(\cdot+h)-f\|_p$。若两个平移量相差大于支集直径，则支集不交，$\|f_n-f_m\|_p^p=2\|f\|_p^p$。任意无限子列都含任意远的两项，故不可能是 Cauchy 列。
\end{solution}

\Needspace{10\baselineskip}
\section{紧算子与核算子}\label{sec:yosida-10-02}
\noindent\textit{原书起始页：中译本 290；英文版 277。}

\begin{definition}[紧算子]
$T\in\mathcal L(X,Y)$ 称为紧算子，若 $T(B_X)$ 相对紧，记 $T\in\mathcal K(X,Y)$。等价地，任何有界列 $x_n$ 的像 $Tx_n$ 都有范数收敛子列。
\end{definition}

\begin{proposition}[闭理想性质]\label{prop:y10-ideal}
紧算子的线性组合仍紧；与有界算子在任意一侧复合仍紧。若紧算子 $T_n$ 在算子范数下收敛于 $T$，则 $T$ 紧。特别地，有限秩算子及其范数极限紧。
\end{proposition}
\begin{proof}
线性组合的像包含在两个相对紧集的和中；右复合将单位球送入一个有界集，左复合则保持紧闭包的连续像紧。给定 $\varepsilon>0$，选 $n$ 使 $\|T-T_n\|<\varepsilon$，再取 $T_n(B_X)$ 的有限 $\varepsilon$ 网，便得到 $T(B_X)$ 的有限 $2\varepsilon$ 网。有限维有界集全有界，故有限秩情形成立。
\end{proof}

本命题只推出有限秩范数极限必紧；一般 Banach 空间上，反方向涉及逼近性质，不在此作无条件断言。

\begin{example}[连续核与平方可积核]
若 $k\in C([a,b]^2)$，则 $Tf(x)=\int_a^b k(x,y)f(y)\,\dd y$ 在 $C[a,b]$ 上紧。事实上 $\|Tf\|_\infty\le(b-a)\|k\|_\infty\|f\|_\infty$，且对 $\|f\|_\infty\le1$，
\[
 |Tf(x)-Tf(z)|\le(b-a)\sup_y|k(x,y)-k(z,y)|.
\]
核的一致连续性和 Ascoli 定理给出紧性。

若 $(S,\mu)$ 为 $\sigma$-有限测度空间，$k\in L^2(S\times S)$，则对应积分算子在 $L^2(S)$ 上紧，且 $\|T\|\le\|k\|_2$。Cauchy--Schwarz 后积分给出此界。矩形乘积函数的有限线性组合在 $L^2(S\times S)$ 中稠密：可先用 $\sigma$-有限分解截到有限测度，再用生成乘积 $\sigma$-代数的矩形逼近简单函数。它们的积分算子为有限秩，且核的 $L^2$ 逼近给出算子范数逼近。这类算子称 Hilbert--Schmidt 积分算子。
\end{example}

\begin{definition}[核算子与核范数]\label{def:y10-nuclear}
若存在 $f_n\in X'$、$y_n\in Y$ 满足
\[
 Tx=\sum_{n=1}^\infty f_n(x)y_n,\qquad
 \sum_{n=1}^\infty\|f_n\|\,\|y_n\|<\infty,
\]
则 $T$ 为核算子。其核范数 $\nu(T)$ 为所有此类表示的上式总和之下确界。这等价于原书使用的 $\ell^1$ 系数乘以两个有界序列的表示。
\end{definition}
\begin{proposition}[核算子必紧]
上式在单位球上一致绝对收敛，部分和 $T_N$ 为有限秩，且
\[
 \|T-T_N\|\le\sum_{n>N}\|f_n\|\,\|y_n\|.
\]
因此核算子紧，$\|T\|\le\nu(T)$，且 $\nu(ATB)\le\|A\|\nu(T)\|B\|$。
\end{proposition}
\begin{proof}
对 $\|x\|\le1$ 逐项估计即得尾和界，紧性由命题 \ref{prop:y10-ideal}。对复合把表示改成 $(f_n\circ B)(x)Ay_n$，再取下确界。
\end{proof}

\begin{proposition}[Sobolev 阶差足够大时的核嵌入]\label{prop:y10-nuclear-sobolev}
设 $G\subset\R^d$ 有界开，$k>j\ge0$ 为整数且 $k-j>d$。则自然嵌入 $H_0^k(G)\to H_0^j(G)$ 为核算子。
\end{proposition}
\begin{proof}
把 $\overline G$ 放在一个周期立方体 $Q$ 的内部。$H_0^s(G)$ 的零延拓是 $H^s(Q)$ 中的闭子空间，因其由测试函数的闭包定义；导数范数与 Fourier 范数等价。周期空间的正交模 $e_m$，$m\in\mathbb Z^d$，在 $H^s(Q)$ 的范数等价于 $(1+|m|^2)^{s/2}$。于是嵌入 $H^k(Q)\to H^j(Q)$ 有秩一展开，归一化后的系数为
\[
 (1+|m|^2)^{-(k-j)/2},\qquad
 \sum_{m\in\mathbb Z^d}(1+|m|^2)^{-(k-j)/2}<\infty.
\]
最后一个和由格点球壳计数 $O(r^{d-1})$ 收敛。将零延拓、此嵌入、$H^j(Q)$ 到零延拓闭子空间的正交投影以及逆零延拓复合，得到所求映射。等价范数不改变核性。
\end{proof}

原书随后联系迹类。这里可直接证明：两个 Hilbert--Schmidt 算子 $A:H\to K$、$B:K\to H$ 的乘积 $BA$ 为核算子。取 $K$ 的正交基 $e_n$，有
\[
 BAx=\sum_n(x,A^*e_n)Be_n,\quad
 \sum_n\|A^*e_n\|\,\|Be_n\|\le\|A\|_{\rm HS}\|B\|_{\rm HS}.
\]
一般迹类与奇异值理论留待谱表示以后讨论。

\begin{exercise}[对角紧算子]\label{ex:y10-diagonal}
设 $D_a:\ell^2\to\ell^2$，$D_ax=(a_nx_n)$，$a\in\ell^\infty$。证明 $D_a$ 紧当且仅当 $a_n\to0$，并证明 $a\in\ell^1$ 时它为核算子。
\end{exercise}
\begin{solution}
截断前 $N$ 项的有限秩算子与 $D_a$ 的范数差为 $\sup_{n>N}|a_n|$。若 $a_n$ 不趋于零，存在子列 $|a_{n_k}|\ge\varepsilon$，其像 $a_{n_k}e_{n_k}$ 两两距离至少 $\sqrt2\varepsilon$，故不紧。核表示是 $D_ax=\sum_n a_nx_ne_n$，绝对范数和为 $\sum_n|a_n|$。
\end{solution}

\begin{exercise}[紧算子改善弱收敛]\label{ex:y10-weak-strong}
若 $T$ 紧且 $x_n\rightharpoonup x$，证明 $Tx_n\to Tx$。
\end{exercise}
\begin{solution}
弱收敛列由一致有界原理有界。若范数收敛失败，可选子列距 $Tx$ 至少 $\varepsilon$；紧性再给出范数收敛子列。它同时弱收敛于 $Tx$，因 $f(Tx_n)=T'f(x_n)$，故范数极限只能为 $Tx$，矛盾。
\end{solution}

\Needspace{10\baselineskip}
\section{Rellich--Gårding 定理}\label{sec:yosida-10-03}
\noindent\textit{原书起始页：中译本 294；英文版 281。}

\begin{theorem}[零边界 Sobolev 紧嵌入]\label{thm:y10-rellich}
设 $G\subset\R^d$ 为有界开集，$k>j\ge0$ 为整数。则 $H_0^k(G)\hookrightarrow H_0^j(G)$ 紧。更一般地，若
\[
 T\in\mathcal L(H_0^k(G),H_0^k(G)),\qquad
 \|T\varphi\|_{H^k}\le C\|\varphi\|_{H^j}
 \quad(\varphi\in C_c^\infty(G)),
\]
则 $T$ 在 $H_0^k(G)$ 上紧。
\end{theorem}
\begin{proof}
令 $u_n$ 在 $H_0^k(G)$ 中有界。由测试函数闭包定义，零延拓 $\widetilde u_n\in H^k(\R^d)$，其导数也是相应导数的零延拓，且范数保持。它们具有共同有界支集，因而 $L^2$ 范数有界且 $L^1$ 范数一致有界。

Fourier 变换满足 $|\widehat {\widetilde u_n}(\xi)|\le C$，且
\[
 |\widehat {\widetilde u_n}(\xi)-\widehat {\widetilde u_n}(\eta)|
 \le |\xi-\eta|\int|x|\,|\widetilde u_n(x)|\,\dd x\le C_1|\xi-\eta|.
\]
在频率闭球上应用 Ascoli，再取对角子列，得到每个频率球上的一致收敛。由 Plancherel 及整数 Sobolev 范数的 Fourier 等价式，频率高于 $R$ 的差满足
\[
 \int_{|\xi|>R}(1+|\xi|^2)^j
 |\widehat {\widetilde u_n}-\widehat {\widetilde u_m}|^2\,\dd\xi
 \le C_2(1+R^2)^{j-k}.
\]
先令 $R$ 大，再令 $m,n$ 大，得到该子列在 $H^j(\R^d)$ 中 Cauchy。限制回 $G$ 后在 $H^j(G)$ 中收敛；每个 $u_n\in H_0^j(G)$，而后者闭，极限仍属于它。这证明嵌入紧。

对第二项，先由 $C_c^\infty$ 在 $H_0^k$ 中稠密，将不等式延至全空间。对任意有界列选出在 $H_0^j$ 中 Cauchy 的子列，不等式使其 $T$ 像在 $H_0^k$ 中 Cauchy，即得紧性。
\end{proof}

这里不要求 $G$ 的边界光滑，因为使用的是零边界闭包空间。若改成 $H^k(G)$，则需另行提供延拓条件，例如有界 Lipschitz 域。全空间上的同类嵌入通常不紧：固定函数的远距离平移仍保持所有 Sobolev 范数。

\begin{example}[正则性如何产生紧性]
设 $S:L^2(G)\to H_0^1(G)$ 有界，则作为 $L^2(G)\to L^2(G)$ 的算子 $JS$ 紧，其中 $J$ 为自然嵌入。紧性来自解算子与紧嵌入的复合；不要求 $S$ 在目标 $H_0^1$ 范数下紧。
\end{example}

\begin{exercise}[平移估计的另一条证明]\label{ex:y10-translation-estimate}
证明 $u\in H^1(\R^d)$ 时，$\|u(\cdot+h)-u\|_2\le |h|\|\nabla u\|_2$。用它及定理 \ref{thm:y10-kolmogorov} 重证 $H_0^1(G)\hookrightarrow L^2(G)$ 紧。
\end{exercise}
\begin{solution}
对光滑函数写差为 $\int_0^1h\cdot\nabla u(x+th)\,\dd t$，由 Minkowski 与平移等距性得界，再由稠密性推广。零延拓的有界族满足此一致估计，且共同有界支集给出尾部条件，故相对紧。
\end{solution}

\begin{exercise}[紧性与核性的阶差]\label{ex:y10-orders}
比较 $H_0^k(G)\to H_0^j(G)$ 的本节结论与命题 \ref{prop:y10-nuclear-sobolev}；为何不应将两者的条件混用？
\end{exercise}
\begin{solution}
紧性只需 $k-j>0$，Fourier 高频因子即趋于零。所给核性充分条件需 $k-j>d$，因为要求格点上的系数绝对可和。系数趋于零与系数可和是不同要求；核性证明提供更强结论，也需要更强估计。
\end{solution}

\Needspace{10\baselineskip}
\section{Schauder 定理}\label{sec:yosida-10-04}
\noindent\textit{原书起始页：中译本 295；英文版 282。}

\begin{theorem}[紧性的对偶刻画]\label{thm:y10-schauder}
有界线性算子 $T:X\to Y$ 紧，当且仅当 $T':Y'\to X'$ 紧。
\end{theorem}
\begin{proof}
若 $T$ 紧，令 $K=\overline{T(B_X)}$，它为紧度量空间。对任意 $f_n\in B_{Y'}$，将其限制于 $K$，得到函数 $F_n(y)=f_n(y)$。它们一致有界且满足
\[
 |F_n(y)-F_n(z)|\le\|y-z\|.
\]
Ascoli 给出在 $K$ 上一致收敛的子列。因此
\[
 \|T'f_n-T'f_m\|
 =\sup_{x\in B_X}|(f_n-f_m)(Tx)|\longrightarrow0.
\]
$X'$ 完备，故 $T'$ 紧。

反之，已证方向应用于 $T'$，得到 $T'':X''\to Y''$ 紧。自然等距嵌入 $J_X,J_Y$ 满足 $T''J_X=J_YT$。故 $J_YT(B_X)$ 相对紧于 $Y''$。$J_Y(Y)$ 是闭子空间，因为 $Y$ 完备；其闭包因此仍在 $J_Y(Y)$ 内。由等距性拉回，$T(B_X)$ 相对紧于 $Y$。
\end{proof}

此证明使用自然嵌入的等距性与闭像，不要求 $X$ 或 $Y$ 自反。Hilbert 空间中经 Riesz 对应可写成 $T$ 紧当且仅当 $T^*$ 紧，但 Banach 空间中的 $T'$ 始终作用在对偶空间上。

\begin{exercise}[对偶方向的估计]\label{ex:y10-adjoint-net}
若 $\{y_1,\ldots,y_N\}$ 是 $T(B_X)$ 的 $\varepsilon$ 网，证明对 $f,g\in B_{Y'}$ 有
\[
 \|T'f-T'g\|\le\max_{1\le i\le N}|(f-g)(y_i)|+2\varepsilon.
\]
\end{exercise}
\begin{solution}
对每个 $x\in B_X$ 选 $i$ 使 $\|Tx-y_i\|<\varepsilon$。因 $\|f-g\|\le2$，有 $|(f-g)(Tx)|\le|(f-g)(y_i)|+2\varepsilon$，对 $x$ 取上确界即得。
\end{solution}

\begin{exercise}[恒等算子]\label{ex:y10-identity}
用 Schauder 定理说明：$I_X$ 紧当且仅当 $I_{X'}$ 紧，并给出这两个条件的维数含义。
\end{exercise}
\begin{solution}
$I_X'=I_{X'}$，故二者等价。由 Riesz 引理，无限维赋范空间的单位球不全有界，所以 $I_X$ 紧当且仅当 $X$ 有限维；这时 $X'$ 也有限维。逆向也可由自然等距嵌入 $X\hookrightarrow X''$ 得到。
\end{solution}

\Needspace{10\baselineskip}
\section{Riesz--Schauder 理论}\label{sec:yosida-10-05}
\noindent\textit{原书起始页：中译本 296；英文版 283。}

\begin{lemma}[闭值域与商空间估计]\label{lem:y10-closed-range}
设 $K\in\mathcal K(X)$，$\lambda\ne0$，$A=\lambda I-K$。则 $N=\ker A$ 有限维，$R(A)$ 闭，且存在 $c>0$ 使
\[
 \|Ax\|\ge c\,\operatorname{dist}(x,N)\quad(x\in X).
\]
\end{lemma}
\begin{proof}
在 $N$ 上 $K=\lambda I$，所以 $N$ 的单位球相对紧；$N$ 闭，故由 Riesz 引理它有限维。

若估计不成立，可按距离归一化并减去核向量，选 $x_n$ 使 $\operatorname{dist}(x_n,N)=1$、$\|x_n\|\le2$、$Ax_n\to0$。取 $Kx_n$ 的收敛子列，利用 $\lambda x_n=Ax_n+Kx_n$ 得 $x_n\to x$，且 $Ax=0$，与距离为 $1$ 矛盾。

若 $Ax_n\to y$，估计使 $\operatorname{dist}(x_n,N)$ 有界。选核向量 $z_n$ 使 $x_n-z_n$ 有界，取 $K(x_n-z_n)$ 收敛子列。同一恒等式给出 $x_n-z_n\to x$，于是 $Ax=y$，所以值域闭。
\end{proof}

\begin{theorem}[Fredholm 择一的可逆性部分]\label{thm:y10-alternative}
对紧算子 $K$ 及 $\lambda\ne0$，$A=\lambda I-K$ 单射当且仅当满射，此时有有界逆。因此紧算子的每个非零谱点都是特征值。
\end{theorem}
\begin{proof}
设 $A$ 单射但不满射。令 $R_n=A^nX$，$R_0=X$。因为 $A^n=\lambda^nI-K_n$，其中 $K_n$ 是紧算子的多项式且常数项为零，上一引理说明每个 $R_n$ 闭。所有包含 $R_{n+1}\subset R_n$ 均严格：若某处相等，$A^nx=A^{n+1}y$，由单射性消去 $A^n$ 得 $x=Ay$，将推出满射。

由 Riesz 引理选 $x_n\in R_n$，$\|x_n\|=1$，$\operatorname{dist}(x_n,R_{n+1})>1/2$。若 $m>n$，则
\[
 \lambda^{-1}(Kx_n-Kx_m)=x_n-x_m-\lambda^{-1}(Ax_n-Ax_m),
\]
右侧除 $x_n$ 外均在 $R_{n+1}$，所以 $\|Kx_n-Kx_m\|>|\lambda|/2$，与紧性矛盾。故单射推出满射。

若 $A$ 满射，$A'$ 单射；$K'$ 由 Schauder 定理紧，已证方向使 $A'$ 满射。若 $Ax=0$，则对任意 $f\in X'$ 可写 $f=A'g$，从而 $f(x)=g(Ax)=0$。Hahn--Banach 分离点，故 $x=0$。有界逆由开映射定理得到。
\end{proof}

\begin{theorem}[非零谱的离散性]\label{thm:y10-discrete-spectrum}
紧算子 $K$ 的非零谱点至多可数，每个均孤立且特征空间有限维。对每个 $\varepsilon>0$，满足 $|\lambda|\ge\varepsilon$ 的谱点只有有限个；故非零谱的唯一可能聚点为零。
\end{theorem}
\begin{proof}
有限维特征空间已证。若存在无穷多个不同特征值 $\lambda_n$ 且 $|\lambda_n|\ge\varepsilon$，由 $|\lambda_n|\le\|K\|$ 取子列使 $\lambda_n\to\lambda\ne0$。取对应特征向量 $v_n$，令 $E_n=\operatorname{span}(v_1,\ldots,v_n)$。不同特征值的特征向量线性无关，故可选 $y_n\in E_n$，$\|y_n\|=1$、$\operatorname{dist}(y_n,E_{n-1})>1/2$。

因 $(K-\lambda_nI)E_n\subset E_{n-1}$，对于 $m<n$，$\lambda_n^{-1}Ky_n-\lambda_m^{-1}Ky_m$ 是 $y_n$ 与 $E_{n-1}$ 中向量之和，其范数大于 $1/2$。但 $Ky_n$ 有收敛子列，且 $\lambda_n\to\lambda\ne0$，这迫使 $\lambda_n^{-1}Ky_n$ 也有收敛子列，矛盾。于是每个远离零的区域只有有限谱点，按 $\varepsilon=1/r$ 取并即得可数性与孤立性。
\end{proof}

\begin{theorem}[有限秩谱投影与相容条件]\label{thm:y10-fredholm}
设 $\lambda\ne0$ 是 $K$ 的谱点，$P$ 为其 Riesz 投影。则 $P$ 有限秩，$\lambda$ 是预解式的有限阶极点；对充分大 $r$，
\[
 X=\ker(\lambda I-K)^r\oplus R((\lambda I-K)^r).
\]
此外，对所有 $\lambda\ne0$，
\begin{align*}
 \dim\ker(\lambda I-K)&=\dim\ker(\lambda I-K'),\\
 (\lambda I-K)x=y\text{ 可解}
 &\Longleftrightarrow f(y)=0\quad\text{对所有 }f\in\ker(\lambda I-K').
\end{align*}
有解时全体解为一个特解加上 $\ker(\lambda I-K)$。
\end{theorem}
\begin{proof}
取只包围 $\lambda$ 且不包围零的小圆 $\Gamma$。由 $R(z,K)=z^{-1}I+z^{-1}R(z,K)K$，得到
\[
 P=\frac1{2\pi i}\int_\Gamma z^{-1}R(z,K)K\,\dd z.
\]
被积算子紧，积分是有限和的算子范数极限，故 $P$ 紧。投影在闭空间 $PX$ 上是恒等算子，所以 $PX$ 有限维。定理 \ref{thm:y8-riesz-projection} 和 \ref{thm:y8-pole} 给出极点、幂零部分和直和分解。

令 $E=PX$、$F=(I-P)X$。$A=\lambda I-K$ 在 $F$ 上可逆，在 $E$ 上为有限维算子。因此 $\dim\ker A=\operatorname{codim}R(A)$，后者由闭值域的湮灭子对应等于 $\dim\ker A'$。若 $\lambda\in\rho(K)$，等式两侧同为零。又 Hahn--Banach 给出闭子空间 $R(A)={}^{\perp}\ker A'$，这恰是相容条件。解集描述由线性性直接得到。
\end{proof}

这里的维数等式是普通特征空间的维数相等；$\operatorname{rank}P$ 则是代数重数，可能严格大于普通特征空间维数。对 Hilbert 伴随，相容条件中的核写成 $\ker(\overline\lambda I-K^*)$，共轭来自 Riesz 对应。

\begin{proposition}[幂紧推广]\label{prop:y10-power-compact}
若 $V\in\mathcal L(X)$ 且某个 $V^q$ 紧，则上述非零谱结构、闭值域与 Fredholm 相容条件仍对 $V$ 成立。
\end{proposition}
\begin{proof}
对 $\lambda\ne0$，若 $\lambda^q\notin\sigma(V^q)$，由
\[
 (\lambda I-V)(\lambda^{q-1}I+\lambda^{q-2}V+\cdots+V^{q-1})
 =\lambda^qI-V^q
\]
及各因子交换可知 $\lambda I-V$ 可逆。若 $\lambda^q\in\sigma(V^q)$，取 $V^q$ 在该谱点的有限秩投影 $Q$。它与 $V$ 交换。$QX$ 有限维，而在 $(I-Q)X$ 上 $\lambda^qI-V^q$ 可逆，故 $\lambda I-V$ 也可逆。问题完全归结为 $QX$ 上的矩阵。因此非零谱均由有限维块产生，闭值域、零指数与极点性质随之成立。远离零的谱点有限，由其 $q$ 次幂的有限性得到。
\end{proof}

\begin{exercise}[共振的相容条件]\label{ex:y10-resonance}
在 $\ell^2$ 上取 $Ke_n=n^{-1}e_n$。解 $(I-K)x=y$ 并给出可解条件及全体解。
\end{exercise}
\begin{solution}
第一坐标要求 $y_1=0$；对 $n\ge2$，$x_n=y_n/(1-1/n)$，乘子至多为 $2$，故所得尾部属于 $\ell^2$。$x_1$ 任意。这正是 $y\perp\ker(I-K)$ 的条件，其中 $\ker(I-K)=\operatorname{span}\{e_1\}$。
\end{solution}

\begin{exercise}[零谱点的例外]\label{ex:y10-zero}
证明无限维 Banach 空间上的紧算子必有 $0\in\sigma(K)$，并说明零不一定是特征值。
\end{exercise}
\begin{solution}
若 $K$ 有有界逆，则 $I=K^{-1}K$ 紧，迫使空间有限维。$\ell^2$ 上 $Ke_n=n^{-1}e_n$ 单射，故其零谱点不是特征值；同理第八章的 Volterra 算子给出另一个例子。
\end{solution}

\begin{exercise}[代数重数与几何重数]\label{ex:y10-jordan}
取 $K=\left(\begin{smallmatrix}a&1\\0&a\end{smallmatrix}\right)\oplus0$ 作用于 $\mathbb C^2\oplus\ell^2$，$a\ne0$。求 $a$ 的 Riesz 投影、两种重数和预解式极点阶数。
\end{exercise}
\begin{solution}
$K$ 有限秩，故紧。$P=I_{\mathbb C^2}\oplus0$，代数重数为 $2$，几何重数为 $1$。左上块逆为 $(z-a)^{-1}I_2+(z-a)^{-2}E_{12}$，故是二阶极点；无限维补空间在 $a$ 附近的逆 $z^{-1}I$ 全纯。
\end{solution}

\Needspace{10\baselineskip}
\section{Dirichlet 问题}\label{sec:yosida-10-06}
\noindent\textit{原书起始页：中译本 299；英文版 286。}

原书用高阶强椭圆算子说明紧性与可解性的联系。这里先在变分形式中给出完整论证，再说明它与分布方程的对应。设 $G\subset\R^d$ 有界开，$V=H_0^m(G)$，$H=L^2(G)$，$J:V\to H$ 为自然嵌入。令 $V^\dagger$ 表示连续共轭线性泛函空间，使 $u\mapsto a(u,\cdot)$ 线性。

\begin{theorem}[Gårding 形式的 Fredholm 择一]\label{thm:y10-dirichlet}
设 $a:V\times V\to\mathbb C$ 是有界半双线性形式，对第一变量线性，且存在 $c>0$、$\beta\ge0$ 使
\[
 \operatorname{Re}a(u,u)+\beta\|Ju\|_H^2\ge c\|u\|_V^2.
\]
则 $Au=a(u,\cdot):V\to V^\dagger$ 的值域闭，核与余核有限维且维数相同。方程 $a(u,v)=F(v)$ 对所有 $v\in V$ 可解，当且仅当
\[
 F(v)=0\quad\text{对所有满足 }a(w,v)=0\ (w\in V)\text{ 的 }v.
\]
特别地，或者齐次问题有非零解，或者对每个 $F\in V^\dagger$ 都有唯一解，且解连续依赖 $F$。
\end{theorem}
\begin{proof}
取实数 $\gamma>\beta$，令 $b(u,v)=a(u,v)+\gamma(Ju,Jv)_H$。由 Lax--Milgram 定理 \ref{thm:y3-lax-milgram}，$B:u\mapsto b(u,\cdot)$ 是 $V$ 到 $V^\dagger$ 的有界同构。定义 $M:V\to V^\dagger$ 为 $Mu(v)=(Ju,Jv)_H$。因为 $J$ 紧而 $H\to V^\dagger$，$h\mapsto(h,J\cdot)_H$ 有界，$M$ 紧。于是
\[
 A=B(I-C),\qquad C=\gamma B^{-1}M\in\mathcal K(V).
\]
Riesz--Schauder 理论说明 $I-C$ 以及 $A$ 值域闭，核与余核维数相同且有限。

为得到具体相容条件，给 $V$ 任选其标准 Hilbert 内积，写 $a(u,v)=(\mathcal Au,v)_V$，$F(v)=(g,v)_V$。因 $\mathcal A$ 值域闭，$g\in R(\mathcal A)$ 等价于 $g\perp\ker\mathcal A^*$。而 $v\in\ker\mathcal A^*$ 恰指 $a(w,v)=0$ 对所有 $w$ 成立，故得到所写条件。若核为零，则余核为零，$A$ 双射，有界逆定理给出稳定估计。
\end{proof}

\begin{example}[高阶算子与非齐次边界数据]
设
\[
 Lu=\sum_{|\alpha|,|\eta|\le m}(-1)^{|\alpha|}
 \partial^\alpha\bigl(a_{\alpha\eta}(x)\partial^\eta u\bigr),\quad
 a(u,v)=\sum_{\alpha,\eta}\int_G a_{\alpha\eta}
 \partial^\eta u\,\overline{\partial^\alpha v}\,\dd x.
\]
假定系数有界，并明确假定此形式满足上述 Gårding 不等式；第 \ref{sec:yosida-06-08} 节提供相应强椭圆、系数正则条件下的来源。给定 $f\in L^2(G)$ 与一个边界提升 $u_1\in H^m(G)$，寻求 $u-u_1\in V$。写 $u=u_1+w$，则
\[
 a(w,v)=(f,Jv)_H-a(u_1,v)=F(v),\qquad v\in V.
\]
右侧属于 $V^\dagger$，故上一定理给出精确可解条件。在形式对称时，相容条件直接对所有齐次解 $v$ 检验；若齐次核为零，每对 $f,u_1$ 都产生唯一弱解，并有
\[
 \|w\|_{H^m}\le C\bigl(\|f\|_2+\|u_1\|_{H^m}\bigr).
\]
测试 $v\in C_c^\infty(G)$ 并分部积分可知这是分布意义的 $Lu=f$。反之，在 $u-u_1\in H_0^m(G)$ 的前提下，分布等式由测试函数稠密性延至 $V$。这里没有把弱解自动提升为 $H^{2m}$ 解；那需要额外正则性与边界假设。任意开集上的 $H_0^m$ 条件也不直接写成逐点边界导数条件。
\end{example}

原书采取 $I+\alpha L$ 的可逆平移，再构造紧算子 $S_1T$；本节用 $A+\gamma M$ 的变分同构实现同一机制。两种写法均先取得强制性，再借助紧嵌入撤去平移。非齐次问题只有在齐次核为零时才对所有数据唯一可解。

\begin{exercise}[区间上的共振]\label{ex:y10-dirichlet-resonance}
在区间 $(0,\pi)$ 上考察零边界问题
\[
 -u''-u=f,\qquad u\in H_0^1(0,\pi).
\]
证明它可解当且仅当 $\int_0^\pi f(x)\sin x\,\dd x=0$，并描述全体解。
\end{exercise}
\begin{solution}
形式为 $a(u,v)=\int u'\overline{v'}-\int u\overline v$，加 $2\|u\|_2^2$ 后强制。齐次弱解满足分布方程 $u''=-u$，在一维可积分或由常系数方程得 $u=A\sin x+B\cos x$；零端点使 $B=0$。形式对称，故相容条件正为对 $\sin x$ 正交。有解时任一特解加 $A\sin x$ 为全体解。
\end{solution}

\begin{exercise}[强制情形没有共振]\label{ex:y10-coercive}
设 $a(u,u)=\|\nabla u\|_2^2+\|u\|_2^2$，$u\in H_0^1(G)$。证明每个 $F\in V^\dagger$ 都有唯一解并给出估计。
\end{exercise}
\begin{solution}
用 $\|u\|_V^2=\|u\|_2^2+\|\nabla u\|_2^2$，形式强制常数为 $1$。Lax--Milgram 给出唯一解；取 $v=u$ 得 $\|u\|_V^2\le\|F\|\|u\|_V$，故 $\|u\|_V\le\|F\|$。此时不必再用紧性撤去平移。
\end{solution}

\Needspace{10\baselineskip}
\section*{本章附录：A. Grothendieck 的核空间}
\phantomsection\addcontentsline{toc}{section}{本章附录：A. Grothendieck 的核空间}
\label{app:yosida-10}
\noindent\textit{原书页码：中译本 302--306；英文版 289--293。}

核空间把一个算子的可和分解提升为局部凸拓扑的性质。以下设 $E$ 为 Hausdorff 局部凸空间。对连续半范数 $p$，记 $E_p=E/\ker p$，范数为 $p$，其 Banach 完备化记为 $\widehat E_p$。若 $q\ge p$，则有范数至多为 $1$ 的自然映射 $i_{q,p}:\widehat E_q\to\widehat E_p$；它未必单射。

\begin{definition}[局部凸空间到 Banach 空间的核映射]
连续线性映射 $T:E\to Y$ 称为核映射，若存在等度连续的 $f_n\in E'$、有界的 $y_n\in Y$ 及 $c\in\ell^1$，使
\[
 Tx=\sum_n c_nf_n(x)y_n.
\]
等度连续意味着存在连续半范数 $p$ 使 $|f_n(x)|\le p(x)$。级数在每个 $p$-有界集上一致绝对收敛，且
\[
 \Bigl\|Tx-\sum_{n\le N}c_nf_n(x)y_n\Bigr\|
 \le p(x)\sup_n\|y_n\|\sum_{n>N}|c_n|.
\]
因此 $T$ 连续，并将某个零邻域送入相对紧集。
\end{definition}

\begin{definition}[核空间]\label{def:y10-nuclear-space}
$E$ 称为核空间，若对每个连续半范数 $p$，存在连续半范数 $q\ge p$，使 $i_{q,p}:\widehat E_q\to\widehat E_p$ 为 Banach 空间间的核算子。这与原书用凸平衡零邻域及其 Minkowski 泛函的定义一致。
\end{definition}

\begin{theorem}[到 Banach 空间的映射判别]\label{thm:y10-nuclear-characterization}
$E$ 为核空间，当且仅当从 $E$ 到任意 Banach 空间的每个连续线性映射都为核映射。
\end{theorem}
\begin{proof}
若 $E$ 核且 $T:E\to Y$ 连续，选连续半范数 $p$ 使 $\|Tx\|\le p(x)$。$T$ 因子分解为 $E\to\widehat E_p\to Y$。再选核映射 $i_{q,p}$，便有分解 $E\to\widehat E_q\to\widehat E_p\to Y$，其可和表示拉回 $E$ 后是核表示。

反之，对自然映射 $j_p:E\to\widehat E_p$ 使用假设，写 $j_px=\sum c_nf_n(x)y_n$。取连续半范数 $q_0$ 使 $|f_n(x)|\le q_0(x)$，令 $q=p+q_0$。每个 $f_n$ 都在 $\ker q$ 上为零，故延为 $\widehat E_q$ 上范数至多为 $1$ 的泛函。级数因此定义核算子 $\widehat E_q\to\widehat E_p$，它在稠密的 $E/\ker q$ 上恰是自然映射，故等于 $i_{q,p}$。
\end{proof}

\begin{proposition}[Hilbert 半范数的来源]\label{prop:y10-hilbert-seminorm}
若 $E$ 核，对任意连续半范数 $p$，都存在由内积诱导的连续半范数 $r\ge p$。因而核空间的拓扑可由 Hilbert 半范数生成。
\end{proposition}
\begin{proof}
将自然映射写成 $j_px=\sum c_nf_n(x)y_n$，吸收相位后可取 $c_n\ge0$，并设 $|f_n(x)|\le q(x)$。定义
\[
 \alpha x=(\sqrt{c_n}f_n(x))_n\in\ell^2,\qquad
 \beta t=\sum_n\sqrt{c_n}t_ny_n\in\widehat E_p.
\]
Cauchy--Schwarz 给出 $\|\alpha x\|_2\le(\sum c_n)^{1/2}q(x)$，且 $\|\beta\|\le(\sum c_n\|y_n\|^2)^{1/2}$。于是 $j_p=\beta\alpha$，从而 $p(x)\le\|\beta\|\|\alpha x\|_2$。令 $r(x)=\max(1,\|\beta\|)\|\alpha x\|_2$ 即可。商掉 $\ker r$ 后是预 Hilbert 空间，其完备化为 Hilbert 空间。
\end{proof}

\begin{example}[乘积空间与快速下降序列]
对任意指标集 $I$，赋乘积拓扑的 $\mathbb C^I$ 为核空间。每个连续半范数 $p$ 都满足 $p(x)\le C\max_{i\in F}|x_i|$，其中 $F$ 有限：这是因为一个控制 $p$ 的基本零邻域仅限制有限坐标，其他坐标可任意伸缩。选右侧为 $q$，则 $\widehat E_q$ 有限维，$i_{q,p}$ 有限秩，因而核。

令 $s$ 为快速下降序列空间，使用 Hilbert 范数
\[
 p_k(x)^2=\sum_{n\ge1}(1+n)^{2k}|x_n|^2,\qquad k=0,1,\ldots.
\]
这些范数与通常的加权上确界半范数定义相同拓扑：上确界由同阶 $p_k$ 控制，而 $p_k$ 由更高一阶的加权上确界控制。$\widehat s_{p_k}$ 是加权 $\ell^2$，其间自然嵌入 $p_{k+2}\to p_k$ 的归一化对角系数为 $(1+n)^{-2}$，绝对可和。任意连续半范数都被某个 $Cp_k$ 控制，再与该核嵌入复合，得到 $s$ 为核空间。
\end{example}

\begin{proposition}[Banach 核空间只能有限维]
若 Banach 空间 $E$ 带其范数拓扑且为核空间，则 $\dim E<\infty$。
\end{proposition}
\begin{proof}
由定理 \ref{thm:y10-nuclear-characterization}，恒等映射 $I_E$ 核，因而紧，Riesz 引理迫使 $E$ 有限维。反之有限维恒等映射有有限秩表示，所以有限维空间核。
\end{proof}

原书还列出核性的子空间、Hausdorff 商、乘积及适当归纳极限稳定性，进而联系 $\Test$、$\Smooth$、$\Sch$ 及其强对偶的核性。这些一般稳定性定理的证明本附录未展开，也不从上述两个例子直接推出。原书末尾的 Minlos 定理把核空间上的适当连续柱分布延拓为测度；仅有有限维相容性并不足以省去连续性条件。本讲义不在此陈述其全部版本或证明，后续专题如需使用应另列准确假设。

\begin{exercise}[等度连续不可省略]\label{ex:y10-equicontinuity}
在 $E=\mathbb C^{\mathbb N}$ 的乘积拓扑上，坐标泛函 $f_n(x)=x_n$ 各自连续。证明 $\{f_n\}$ 不等度连续。
\end{exercise}
\begin{solution}
每个零邻域都包含只限制有限坐标的基本零邻域。任选未受限制的坐标 $n$，该坐标可任意大，故不存在零邻域使所有 $|f_n(x)|\le1$。因此逐个连续不等于等度连续；甚至 $\sum2^{-n}x_n$ 也不对所有乘积空间元素收敛。
\end{solution}

\begin{exercise}[自然映射不必单射]\label{ex:y10-canonical-kernel}
在 $E=\mathbb C^2$ 上取 $q(x)=|x_1|+|x_2|$、$p(x)=|x_1|$。计算 $i_{q,p}$ 的核，并说明这为何不妨碍核空间定义。
\end{exercise}
\begin{solution}
$\widehat E_q=\mathbb C^2$，$\widehat E_p\cong\mathbb C$，自然映射为 $(x_1,x_2)\mapsto x_1$，核为 $\{0\}\times\mathbb C$。核空间的条件要求映射具有可和秩一展开，并不要求其单射；此例为有限秩映射，当然满足要求。
\end{solution}


\Needspace{6\baselineskip}
\section*{本章参考说明与阅读衔接}
\phantomsection\addcontentsline{toc}{section}{本章参考说明与阅读衔接}
直接来源为吉田第十章及章内附录 \cite{yosida-zh,yosida-en}；原书所引 Grothendieck、Gårding、Minlos 等原始文献未计为直接阅读。讲义补充有限网、商空间估计与变分形式的推导。核空间的一般稳定性和 Minlos 定理仅说明位置，未在此证明。下一章将从有限秩谱投影进一步走向正规与自伴算子的谱表示。
